\documentclass[11pt,a4paper]{article}

\usepackage[utf8]{inputenc}
\usepackage[T1]{fontenc}
\usepackage{lmodern}
\usepackage{amsmath,amssymb,amsthm}
\usepackage{booktabs}
\usepackage{tabularx}
\usepackage{hyperref}
\usepackage[margin=2.5cm]{geometry}
\usepackage{natbib}
\usepackage{enumitem}
\usepackage{float}
\usepackage{xcolor}

\newtheorem{theorem}{Theorem}
\newtheorem{proposition}[theorem]{Proposition}
\newtheorem{corollary}[theorem]{Corollary}
\newtheorem{lemma}[theorem]{Lemma}
\newtheorem{definition}[theorem]{Definition}
\newtheorem{remark}[theorem]{Remark}
\newtheorem{example}[theorem]{Example}
\newtheorem{conjecture}[theorem]{Conjecture}
\newtheorem{openquestion}{Open Question}

\newcolumntype{Y}{>{\raggedright\arraybackslash}X}

\newcommand{\SU}{\mathrm{SU}(2)}
\newcommand{\Real}{\operatorname{Re}}
\newcommand{\Imag}{\operatorname{Im}}
\newcommand{\Hq}{\mathbb{H}}

\hypersetup{
  colorlinks=true,
  linkcolor=blue!60!black,
  citecolor=green!40!black,
  urlcolor=blue!60!black
}

\title{Toward a Non-Commutative Residuated Lattice\\
from Quaternion Multiplication}

\author{J.~Arturo Ornelas Brand\\
\small Independent Researcher, Aguascalientes, M\'exico\\
\small \texttt{arturoornelas62@gmail.com}}

\date{Draft --- \today}

\begin{document}

\maketitle

\begin{abstract}
The quaternionic G-lattice introduced in the companion paper uses
component-wise min as its monoidal operation, making the residuated
lattice a product of four independent chains. The quaternionic structure
of the framework manifests in the $\SU$ group action and the product on
$S^3$, but not in the logic's conjunction. This paper investigates
whether the Hamilton product of quaternions can serve as a genuinely
non-commutative conjunction---a monoidal operation that couples the four
epistemic axes through the quaternionic cross terms.

We establish positive results: the Hamilton product on
$V = [0,1] \times [-1,1]^3$ satisfies bilateral identity, bilateral
annihilation, and associativity; it defines left and right implications
via quaternionic division, both satisfying modus ponens; and under every
commutative restriction (Boolean, fuzzy, modal), the two implications
coincide and the product reduces to the standard commutative t-norm.
Non-commutativity is proven to be \emph{emergent}: it arises only in the
full four-dimensional space, not in any proper restriction.

We also characterize fundamental obstructions. The Hamilton product
is not closed on $V$ and is not monotone with respect to the
component-wise order. We prove that no total order on $\Hq$ is
compatible with Hamilton multiplication, that the component-wise
lattice order and the Hamilton monoid cannot coexist on $V$, and---via
four independent impossibility results---that no non-trivial lattice
order on $B^4 = \{q \in \Hq : |q| \leq 1\}$ is compatible with the
Hamilton monoid. The four results are: (i)~no proper bi-invariant cone
exists in $\Hq$ (SU(2) transitivity), (ii)~no non-trivial
bi-invariant partial order exists on $Q_8$ (exhaustive closure),
(iii)~no real-dominant subcone is closed under multiplication
(analytical bound), and (iv)~the norm order collapses the structure to
a one-dimensional product BL-algebra.

These results culminate in a \textbf{quaternionic trilemma}: for any
monoidal operation on a subset of $V = [0,1] \times [-1,1]^3$, at
most two of \{associativity, non-commutativity, compatible lattice
order\} can hold simultaneously. The trilemma is a \emph{theorem}
for $V$-domains (Theorem~\ref{thm:trilemma-V}): De Moivre's theorem
shows that every non-real element's iterates eventually exit $V$, so
the only sub-monoids of $V$ under Hamilton product are commutative.
For the full ball $B^4$, where powers remain bounded but the
angular argument fails, Theorem~\ref{thm:trilemma-B4} proves that
the canonical non-commutative sub-monoid (equal-norm generators)
admits no compatible lattice order. The equal-norm hypothesis is
sharp: for generators of unequal norm, compatible monotone lattice
orders exist (Remark~\ref{rem:sharpness}).

An $\alpha$-parameterized family interpolates between the diagonal
product ($\alpha = 0$) and the full Hamilton product ($\alpha = 1$),
preserving all specializations; associativity holds only at the
endpoints.
\end{abstract}

\noindent\textbf{Keywords:} non-commutative residuated lattice;
quaternion multiplication; t-norm; substructural logic; fuzzy logic;
Hamilton product; non-commutative conjunction

%======================================================================
\section{Introduction}
\label{sec:intro}
%======================================================================

The quaternionic logic developed in \citet{p11} assigns to each
semantic state a four-dimensional coordinate
$(r, i, j, k) \in V = [0,1] \times [-1,1]^3$, where $r$ encodes
crystallised truth, $i$ encodes potentiality, $j$ encodes
verifiability, and $k$ encodes recursive selection. Logical connectives
form a bounded residuated lattice on $V$, and an $\SU$ group acts by
conjugation, yielding a G-lattice whose central result is a commutation
theorem: connectives commute with the group action if and only if they
depend exclusively on the real component.

The lattice connectives---AND~$= \min$, OR~$= \max$, and the G\"odel
residuum---are defined component-wise. The resulting lattice is
isomorphic to the product of four independent chains
\citep[Remark~5]{p11}. The quaternionic structure manifests in the group
action and the product on $S^3$, but not in the conjunction itself.

This is a structural limitation. In the four-dimensional space, the
Hamilton product couples all four components: the relation
$ij = k$ (potentiality composed with verification yields selection) is a
cross-term that no component-wise operation can reproduce. If the
conjunction were the Hamilton product rather than component-wise min, the
logic would be genuinely non-commutative---the order of premises would
matter, as in Lambek calculus \citep{lambek1958} and substructural
logics \citep{restall2000}.

This paper investigates whether such a construction is possible. The
question is: \emph{can a residuated lattice be built on $V$ whose
monoidal operation derives from quaternion multiplication?}

The answer is: partially. We establish the monoid and its residuals,
characterize the obstructions to the full lattice, and map the path
forward.

\paragraph{Contributions.}
\begin{enumerate}[nosep]
\item The Hamilton product on $V$ satisfies bilateral identity,
  bilateral annihilation, and associativity
  (Theorem~\ref{thm:monoid-properties}).
\item Two obstructions are characterized: failure of closure
  (Theorem~\ref{thm:closure-failure}) and failure of monotonicity
  (Theorem~\ref{thm:monotonicity-failure}), both traced to the signed
  structure of the quaternionic cross terms.
\item The unit ball $B^4$ resolves closure
  (Proposition~\ref{prop:ball-monoid}).
\item Left and right implications via quaternionic division satisfy modus
  ponens, tautology, and non-commutativity
  (Theorem~\ref{thm:implications}).
\item Under every commutative restriction (Boolean, fuzzy, modal), the
  two implications coincide and the product reduces to the standard
  t-norm (Theorem~\ref{thm:specialization}).
\item Non-commutativity is emergent: it arises only in the full 4D
  space (Theorem~\ref{thm:emergence}).
\item No total order on $\Hq$ is compatible with Hamilton multiplication
  (Proposition~\ref{prop:no-total-order}).
\item Six independent impossibility results rule out every natural
  candidate for a compatible lattice order on $B^4$: no bi-invariant
  cone (Theorem~\ref{thm:cone-impossibility}), no bi-invariant order
  on $Q_8$ (Theorem~\ref{thm:q8-impossibility}), no closed
  real-dominant subcone (Proposition~\ref{prop:subcone}),
  norm-order collapse to one dimension
  (Proposition~\ref{prop:norm-collapse}), component-wise
  incompatibility (Theorem~\ref{thm:order-incompatibility}), and
  sphere exclusion for integral RLs
  (Proposition~\ref{prop:sphere-exclusion}).
\item The \textbf{quaternionic trilemma}, proved for $V$-domains
  (Theorem~\ref{thm:trilemma-V}) via angular accumulation
  (Theorem~\ref{thm:angular}): at most two of \{associativity,
  non-commutativity, compatible lattice order\} hold simultaneously.
  Proved for the canonical $B^4$ sub-monoid with equal-norm
  generators (Theorem~\ref{thm:trilemma-B4}); the equal-norm
  condition is sharp (Remark~\ref{rem:sharpness}).
\item An $\alpha$-parameterized family bridges the commutative and
  non-commutative regimes, with associativity only at the endpoints
  (Proposition~\ref{prop:alpha}).
\end{enumerate}

%======================================================================
\section{Preliminaries}
\label{sec:prelim}
%======================================================================

\subsection{Quaternion algebra}

The quaternion algebra $\Hq$ is the four-dimensional real algebra
spanned by $\{1, i, j, k\}$ with $i^2 = j^2 = k^2 = ijk = -1$.
Every quaternion $q = a + bi + cj + dk$ has real part $\Real(q) = a$,
imaginary part $\Imag(q) = (b, c, d)$, conjugate
$\bar{q} = a - bi - cj - dk$, and norm
$|q| = \sqrt{a^2 + b^2 + c^2 + d^2}$. The norm is multiplicative:
$|q_1 q_2| = |q_1| \cdot |q_2|$.

\subsection{Hamilton product}

For $q_1 = (r_1, i_1, j_1, k_1)$ and $q_2 = (r_2, i_2, j_2, k_2)$,
the Hamilton product is:
\begin{equation}
\label{eq:hamilton}
q_1 \cdot q_2 = \begin{pmatrix}
  r_1 r_2 - i_1 i_2 - j_1 j_2 - k_1 k_2 \\
  r_1 i_2 + i_1 r_2 + j_1 k_2 - k_1 j_2 \\
  r_1 j_2 - i_1 k_2 + j_1 r_2 + k_1 i_2 \\
  r_1 k_2 + i_1 j_2 - j_1 i_2 + k_1 r_2
\end{pmatrix}
\end{equation}
The product is associative, has identity $1 = (1, 0, 0, 0)$, and every
nonzero quaternion has an inverse $q^{-1} = \bar{q}/|q|^2$.

\subsection{Non-commutative residuated lattices}

A \emph{residuated lattice} \citep{galatos2007} is an algebra
$(L, \wedge, \vee, \cdot, \backslash, /, e)$ where $(L, \wedge, \vee)$
is a lattice, $(L, \cdot, e)$ is a monoid, and the \emph{residuation
law} holds:
\begin{equation}
\label{eq:residuation}
a \cdot b \leq c
  \quad\Longleftrightarrow\quad
  a \leq c / b
  \quad\Longleftrightarrow\quad
  b \leq a \backslash c
\end{equation}
where $\backslash$ is the left residual and $/$ is the right residual.
When $\cdot$ is commutative, the two residuals coincide and we write
$a \to b$.

In the commutative case, the monoidal operation $\cdot$ is a t-norm:
a binary operation on $[0,1]$ that is commutative, associative,
monotone, and has identity~$1$. The G\"odel t-norm ($\min$), the product
t-norm ($a \cdot b$), and the {\L}ukasiewicz t-norm
($\max(0, a + b - 1)$) are the three basic t-norms
\citep{hajek1998,esteva2001}. The companion paper \citet{p11} uses
$\min$ on each component, yielding a product of four G\"odel chains.

Non-commutative extensions were introduced by
\citet{hajek2003} (pseudo-BL algebras) and studied extensively by
\citet{ciungu2006} and \citet{galatos2007}. The key difference from the
commutative case is the existence of \emph{two} residuals, corresponding
to left and right division.

\subsection{Quaternionic fuzzy sets}

\citet{dai2023} introduced quaternion-valued membership functions for
fuzzy sets, defining intersection and union via component-wise min and
max. However, Dai's construction does not define a t-norm, residuum, or
any inference mechanism---it remains at the level of set operations.

%======================================================================
\section{The Hamilton Product as Conjunction}
\label{sec:conjunction}
%======================================================================

We examine which properties of a monoidal operation the Hamilton product
satisfies on $V = [0,1] \times [-1,1]^3$.

\begin{theorem}[Monoid-like properties]
\label{thm:monoid-properties}
The Hamilton product on $V$ satisfies:
\begin{enumerate}[nosep]
\item \textbf{Bilateral identity}: $q \cdot \mathbf{1} = \mathbf{1} \cdot q = q$
  for all $q \in V$, where $\mathbf{1} = (1, 0, 0, 0)$.
\item \textbf{Bilateral annihilation}: $q \cdot \mathbf{0} = \mathbf{0} \cdot q = \mathbf{0}$
  for all $q \in V$, where $\mathbf{0} = (0, 0, 0, 0)$.
\item \textbf{Associativity}: $(q_1 \cdot q_2) \cdot q_3 = q_1 \cdot (q_2 \cdot q_3)$
  for all $q_1, q_2, q_3 \in V$.
\end{enumerate}
\end{theorem}

\begin{proof}
(1)~From \eqref{eq:hamilton} with $q_2 = (1, 0, 0, 0)$: the real-times-real
term gives $r_1$, all cross terms vanish. The calculation is symmetric.

(2)~Every term in \eqref{eq:hamilton} contains a factor from $q_2$; if
$q_2 = \mathbf{0}$, all terms vanish. By symmetry of the bilinear
structure.

(3)~This is the associativity of the quaternion algebra $\Hq$, which
holds for all quaternions, not just those in $V$.
\end{proof}

The Hamilton product also satisfies the basis relations that give it
semantic content:

\begin{proposition}[Basis relations]
\label{prop:basis}
The pure basis quaternions satisfy:
\begin{alignat}{3}
  i \cdot j &= k,  &\qquad j \cdot i &= -k, \label{eq:ij} \\
  j \cdot k &= i,  &\qquad k \cdot j &= -i, \label{eq:jk} \\
  k \cdot i &= j,  &\qquad i \cdot k &= -j, \label{eq:ki} \\
  i^2 &= j^2 = k^2 &&= -\mathbf{1}. \label{eq:squares}
\end{alignat}
\end{proposition}

\begin{remark}[Semantic reading]
The non-commutativity $ij = k \neq ji = -k$ has a direct semantic
interpretation: \emph{imagining first, then verifying} ($ij = k$:
selection) produces a different result from \emph{verifying first, then
imagining} ($ji = -k$: anti-selection). The order of epistemic
operations matters. This is the algebraic signature of the distinction
between hypothesis generation (imagine, then test) and falsification
(test, then reimagine).
\end{remark}

However, two critical properties fail:

\begin{theorem}[Closure failure]
\label{thm:closure-failure}
The Hamilton product is not closed on $V = [0,1] \times [-1,1]^3$.
Specifically, for $q_1, q_2 \in V$ with $\Imag(q_1) \neq 0$ and
$\Imag(q_2) \neq 0$:
\begin{equation}
\Real(q_1 \cdot q_2) = r_1 r_2 - (i_1 i_2 + j_1 j_2 + k_1 k_2)
\end{equation}
which can be negative (violating $r \in [0,1]$) when
$\langle \Imag(q_1), \Imag(q_2) \rangle > r_1 r_2$.
\end{theorem}

\begin{proof}
Take $q_1 = (0.1, 0.8, 0, 0)$ and $q_2 = (0.1, 0.8, 0, 0)$.
Then $\Real(q_1 \cdot q_2) = 0.01 - 0.64 = -0.63 < 0$. Thus
$q_1 \cdot q_2 \notin V$.
\end{proof}

\begin{theorem}[Monotonicity failure]
\label{thm:monotonicity-failure}
The Hamilton product is not monotone with respect to the component-wise
order on $V$. That is, there exist $a, b, c \in V$ with $b \leq c$
(component-wise) but $a \cdot b \not\leq a \cdot c$.
\end{theorem}

\begin{proof}
The signed cross terms in \eqref{eq:hamilton} are the source. Consider
the real component of $a \cdot b$:
\[
(a \cdot b)_r = a_r b_r - a_i b_i - a_j b_j - a_k b_k
\]
Increasing $b_i$ while holding all else fixed \emph{decreases} the
real component of the product (due to the $-a_i b_i$ term). But in
the component-wise order, increasing $b_i$ means $b' \geq b$. So
$(a \cdot b')_r < (a \cdot b)_r$ while $b' \geq b$: left-monotonicity
fails on the real component.

Explicitly: let $a = (0.5, 0.5, 0, 0)$, $b = (0.5, 0.1, 0, 0)$,
$c = (0.5, 0.9, 0, 0)$. Then $b \leq c$ component-wise, but
$(a \cdot b)_r = 0.25 - 0.05 = 0.20$ while
$(a \cdot c)_r = 0.25 - 0.45 = -0.20$. So $a \cdot b > a \cdot c$
on the real component.
\end{proof}

\begin{remark}[The fundamental obstruction]
\label{rem:obstruction}
Both failures---closure and monotonicity---trace to the same algebraic
source: the \textbf{signed cross terms} in the Hamilton product. The
real component formula
$r_1 r_2 - \langle \Imag(q_1), \Imag(q_2) \rangle$ contains a
subtraction of the imaginary dot product. In a standard (commutative)
t-norm on $[0,1]$, all terms are positive, ensuring closure and
monotonicity. In the Hamilton product, the cross terms introduce
\emph{negative contributions} that can dominate the positive diagonal
terms. This is not a defect---it is the algebraic mechanism that
produces non-commutativity and axis coupling ($ij = k$)---but it is
incompatible with a positively-ordered bounded domain.
\end{remark}

%======================================================================
\section{The Unit Ball Monoid}
\label{sec:ball}
%======================================================================

The closure failure on $V$ can be resolved by working on the unit ball.

\begin{proposition}[Unit ball monoid]
\label{prop:ball-monoid}
The set $B^4 = \{q \in \Hq : |q| \leq 1\}$ with the Hamilton product
forms a monoid:
\begin{enumerate}[nosep]
\item \textbf{Closure}: If $|q_1| \leq 1$ and $|q_2| \leq 1$, then
  $|q_1 \cdot q_2| = |q_1| \cdot |q_2| \leq 1$.
\item \textbf{Associativity}: Inherited from $\Hq$.
\item \textbf{Identity}: $\mathbf{1} = (1, 0, 0, 0)$ with $|\mathbf{1}| = 1 \in B^4$.
\end{enumerate}
\end{proposition}

\begin{proof}
Closure follows from the multiplicativity of the quaternion norm.
\end{proof}

\begin{remark}[Relation to $V$]
The domain $V = [0,1] \times [-1,1]^3$ and the unit ball $B^4$ are
different subsets of $\Hq$. Their intersection
$V \cap B^4 = \{q \in V : |q| \leq 1\}$ is nonempty and contains
$\mathbf{0}$ and $\mathbf{1}$. The unit ball $B^4$ is closed under
the Hamilton product; $V$ is not. However, $B^4$ is not closed under
the lattice operations $\min$ and $\max$ (component-wise max of two
elements in $B^4$ can leave $B^4$), so the lattice and monoid
structures live on incompatible domains. This domain incompatibility
is the geometric face of the algebraic obstruction identified in
Remark~\ref{rem:obstruction}.
\end{remark}

The unit ball monoid admits left and right division:

\begin{proposition}[Division on $B^4$]
\label{prop:division}
For $a \in B^4$ with $a \neq \mathbf{0}$:
\begin{align}
  a^{-1} &= \bar{a}\,/\,|a|^2 \label{eq:inv} \\
  a \backslash b &= a^{-1} \cdot b \quad\text{(left division)} \label{eq:ldiv} \\
  b / a &= b \cdot a^{-1} \quad\text{(right division)} \label{eq:rdiv}
\end{align}
Left division satisfies $a \cdot (a \backslash b) = b$ and right
division satisfies $(b / a) \cdot a = b$.
\end{proposition}

\begin{proof}
By associativity: $a \cdot (a^{-1} \cdot b) = (a \cdot a^{-1}) \cdot b = b$.
Similarly for the right division.
\end{proof}

\begin{remark}[Division leaves $B^4$]
\label{rem:div-leaves}
For $a, b \in B^4$ with $|a| < |b|$, the division $a \backslash b$ has
norm $|b|/|a| > 1$, so $a \backslash b \notin B^4$. This means the
residuals of the Hamilton product, defined as suprema via the
residuation law~\eqref{eq:residuation}, do not stay within $B^4$ in
general. A truncation $\min(1, |b|/|a|)$ applied to the norm is the
natural remedy---analogous to the $\min(1, b/a)$ in the product
t-norm's residuum on $[0,1]$---but it breaks the exact adjunction
property. Remark~\ref{rem:residuation-norm} characterises the
scope of this limitation.
\end{remark}

%======================================================================
\section{Two Implications from Non-Commutativity}
\label{sec:implications}
%======================================================================

The non-commutativity of the Hamilton product gives rise to two distinct
implication operators, corresponding to left and right division. We
study their properties on $\Hq \setminus \{0\}$ (the natural domain
where division is defined).

\begin{definition}[Left and right implications]
\label{def:implications}
For $a, b \in \Hq$ with $a \neq \mathbf{0}$:
\begin{align}
  a \to_L b &= a^{-1} \cdot b \quad\text{(left implication)} \\
  a \to_R b &= b \cdot a^{-1} \quad\text{(right implication)}
\end{align}
\end{definition}

\begin{theorem}[Properties of quaternionic implications]
\label{thm:implications}
Both implications satisfy:
\begin{enumerate}[nosep]
\item \textbf{Tautology}: $a \to_L a = a \to_R a = \mathbf{1}$ for all
  $a \neq \mathbf{0}$.
\item \textbf{Modus ponens (left)}: $a \cdot (a \to_L b) = b$.
\item \textbf{Modus ponens (right)}: $(a \to_R b) \cdot a = b$.
\item \textbf{Non-commutativity}: $a \to_L b \neq a \to_R b$ in general.
\end{enumerate}
\end{theorem}

\begin{proof}
(1)~$a^{-1} \cdot a = \mathbf{1}$ and $a \cdot a^{-1} = \mathbf{1}$.

(2)~$a \cdot (a^{-1} \cdot b) = (a a^{-1}) \cdot b = b$ by
associativity.

(3)~$(b \cdot a^{-1}) \cdot a = b \cdot (a^{-1} a) = b$ by
associativity.

(4)~$a^{-1} b \neq b a^{-1}$ whenever $a^{-1}$ and $b$ do not commute.
For example, $i \to_L j = i^{-1} \cdot j = (-i) \cdot j = -k$,
while $i \to_R j = j \cdot i^{-1} = j \cdot (-i) = k$. The left
and right implications differ by sign on the $k$ axis.
\end{proof}

\begin{remark}[Semantic reading of the two implications]
The left implication $a \to_L b = a^{-1} b$ asks: ``given $a$ as
premise, what must I compose on the right to reach $b$?'' The right
implication $a \to_R b = b a^{-1}$ asks: ``what must I compose on the
left to reach $b$, given $a$ as premise?'' In epistemic terms: the left
implication assumes the premise acts first (and asks what follows); the
right implication assumes the conclusion is the target (and asks what
precedes). The two coincide when the order of operations does not
matter---i.e., when the product is commutative.
\end{remark}

\begin{proposition}[Contraposition]
\label{prop:contraposition}
Classical contraposition ($a \to b \iff \neg b \to \neg a$) does
\emph{not} hold for the quaternionic implications. With conjugation as
negation ($\neg q = \bar{q}$):
\begin{equation}
a \to_L b = a^{-1} b
  \quad\neq\quad
  \bar{b} \to_L \bar{a} = \bar{b}^{-1} \bar{a}
\end{equation}
in general.
\end{proposition}

\begin{proof}
$a^{-1} b = \bar{a} b / |a|^2$ while
$\bar{b}^{-1} \bar{a} = b \bar{a} / |b|^2$. These are equal only when
$|a| = |b|$ and $\bar{a} b = b \bar{a}$, which fails generically.
\end{proof}

%======================================================================
\section{Specialization Recovery}
\label{sec:specialization}
%======================================================================

A construction on the full four-dimensional space must reduce to known
logics under the appropriate restrictions. We verify this for the
Hamilton product.

\begin{theorem}[Specialization and commutativity]
\label{thm:specialization}
Under each of the following restrictions, the Hamilton product becomes
commutative and reduces to a standard t-norm, and the two implications
coincide:
\begin{enumerate}[nosep]
\item \textbf{Boolean} ($i = j = k = 0$, $r \in \{0, 1\}$):
  The product is $r_1 \cdot r_2$ (multiplication in $\{0,1\}$), which
  is Boolean AND.
\item \textbf{Fuzzy} ($i = j = k = 0$, $r \in [0,1]$):
  The product is $r_1 \cdot r_2$ (the product t-norm
  $T_{\mathrm{prod}}$).
\item \textbf{Modal} ($j = k = 0$, $r, i \in [0,1]$):
  The product is $(r_1 r_2 - i_1 i_2, \; r_1 i_2 + i_1 r_2, \; 0,
  \; 0)$, which is complex multiplication restricted to the first
  quadrant. This is commutative.
\end{enumerate}
In each case, the Hamilton cross terms involving $j$ and $k$ vanish,
eliminating the source of non-commutativity.
\end{theorem}

\begin{proof}
(1)~With $i_1 = j_1 = k_1 = i_2 = j_2 = k_2 = 0$,
equation~\eqref{eq:hamilton} gives $(r_1 r_2, 0, 0, 0)$.

(2)~Same calculation with $r_1, r_2 \in [0,1]$.

(3)~With $j_1 = k_1 = j_2 = k_2 = 0$, the cross terms
$j_1 k_2 - k_1 j_2 = 0$, $k_1 i_2 - i_1 k_2 = 0$, and
$i_1 j_2 - j_1 i_2 = 0$ all vanish. The remaining product is
$(r_1 r_2 - i_1 i_2, \; r_1 i_2 + i_1 r_2, \; 0, \; 0)$, which is
the multiplication of complex numbers $z_1 = r_1 + i_1 \cdot i$ and
$z_2 = r_2 + i_2 \cdot i$. Complex multiplication is commutative.
\end{proof}

\begin{corollary}[Implications coincide under restriction]
\label{cor:imp-coincide}
Under any of the three restrictions above, $a \to_L b = a \to_R b$ for
all $a, b$ in the restricted domain.
\end{corollary}

\begin{proof}
Commutativity implies $a^{-1} b = b a^{-1}$.
\end{proof}

\begin{remark}[Where non-commutativity begins]
The minimal restriction that produces non-commutativity requires at
least two distinct imaginary axes to be active. Specifically, the
cross terms in~\eqref{eq:hamilton} are:
\begin{equation}
j_1 k_2 - k_1 j_2 \quad\text{(i-component)}, \quad
k_1 i_2 - i_1 k_2 \quad\text{(j-component)}, \quad
i_1 j_2 - j_1 i_2 \quad\text{(k-component)}
\end{equation}
Each cross term couples \emph{two different} imaginary axes. When
only one imaginary axis is active ($i$ only, or $j$ only, or $k$ only),
all cross terms vanish and the product is commutative. Non-commutativity
requires interaction between at least two of $\{i, j, k\}$.
\end{remark}

%======================================================================
\section{Emergent Non-Commutativity}
\label{sec:emergence}
%======================================================================

\begin{theorem}[Non-commutativity is emergent]
\label{thm:emergence}
Let $S \subsetneq \{r, i, j, k\}$ be a proper subset of the four axes
and $V_S = \{q \in V : q_m = 0 \text{ for } m \notin S\}$ the
corresponding restriction. If $|S \cap \{i, j, k\}| \leq 1$, then
the Hamilton product is commutative on $V_S$.
\end{theorem}

\begin{proof}
If the active imaginary axes are $\{i\}$, $\{j\}$, or $\{k\}$ (or
empty), then all cross terms in~\eqref{eq:hamilton} vanish (each
cross term couples two distinct imaginary axes). The product reduces
to real-times-real plus a single imaginary-squared term, which
commutes.

For zero imaginary axes ($S \subseteq \{r\}$): the product is
$r_1 r_2$, trivially commutative.

For one imaginary axis, say $S = \{r, i\}$: the product is
$(r_1 r_2 - i_1 i_2, \; r_1 i_2 + i_1 r_2) = z_1 z_2$
(complex multiplication), which is commutative.
\end{proof}

\begin{corollary}
\label{cor:min-axes}
The Hamilton product is non-commutative on $V_S$ if and only if
$|S \cap \{i, j, k\}| \geq 2$.
\end{corollary}

\begin{proof}
The ``only if'' direction is Theorem~\ref{thm:emergence}. For the
``if'' direction: with two imaginary axes active, at least one cross
term is nonzero. For example, with $i$ and $j$ active:
$q_1 = (0, a, 0, 0)$, $q_2 = (0, 0, b, 0)$ gives
$q_1 q_2 = (0, 0, 0, ab)$ while $q_2 q_1 = (0, 0, 0, -ab)$.
\end{proof}

\begin{remark}[Emergence as dimensional phenomenon]
Non-commutativity is not present in any one- or two-dimensional
restriction of the quaternionic logic. It emerges only when at least
three dimensions are active (real plus two imaginary). This is
analogous to the classical result that the only associative division
algebras over $\mathbb{R}$ are $\mathbb{R}$ (1D, commutative),
$\mathbb{C}$ (2D, commutative), and $\Hq$ (4D, non-commutative): the
non-commutativity is a genuinely four-dimensional phenomenon that cannot
be observed from any lower-dimensional shadow.

In the semantic interpretation: imagining ($i$), verifying ($j$), and
selecting ($k$) are each individually order-independent. It is only
when \emph{two or more} of these processes interact that the order of
operations begins to matter.
\end{remark}

%======================================================================
\section{The Lattice Order Problem}
\label{sec:obstruction}
%======================================================================

The results so far establish a monoid $(B^4, \cdot, \mathbf{1})$ with
left and right residuals on the punctured ball. What is missing for a
residuated lattice is a compatible lattice order.

\begin{proposition}[No compatible total order]
\label{prop:no-total-order}
There is no total order $\leq$ on $\Hq$ such that $(\Hq, \leq)$ is a
totally ordered group under Hamilton multiplication.
\end{proposition}

\begin{proof}
In a totally ordered group, every square is non-negative:
$a^2 \geq e$ (since either $a \geq e$ or $a^{-1} \geq e$, and
$a^2 = a \cdot a \geq a \cdot e = a \geq e$ or
$a^2 = (a^{-1})^{-2} \geq e$). But $i^2 = -1 < 1 = e$ in any
order that makes $e$ the identity. Contradiction.
\end{proof}

\begin{theorem}[Component-wise order incompatibility]
\label{thm:order-incompatibility}
Let $\leq_V$ be the component-wise partial order on $V$. There is no
subset $D \subseteq V$ containing $\mathbf{0}$ and $\mathbf{1}$ such
that $(D, \leq_V, \cdot_{\Hq})$ is a residuated lattice---i.e., such
that $D$ is closed under $\min$, $\max$, and $\cdot_{\Hq}$
simultaneously.
\end{theorem}

\begin{proof}
Suppose such $D$ exists. Since $\mathbf{0}, \mathbf{1} \in D$ and $D$
is closed under $\min$ and $\max$, $D$ contains all elements of the
form $(r, 0, 0, 0)$ for $r \in [0,1]$ (as $\max(\mathbf{0},
\min(\mathbf{1}, (r, 0, 0, 0))) = (r, 0, 0, 0)$, taking
$D$ to contain a generating set).

For $D$ to be a non-trivial quaternionic lattice (not merely the real
line), $D$ must contain some element $a = (r_a, i_a, 0, 0)$ with
$i_a > 0$. Closure under Hamilton product requires
$a^2 = (r_a^2 - i_a^2, \; 2 r_a i_a, \; 0, \; 0) \in D$. If
$i_a > r_a$, then $\Real(a^2) = r_a^2 - i_a^2 < 0$, so
$a^2 \notin V$ (the real component is negative). Thus $D$ cannot
contain any element with $i_a > r_a$.

But even restricting to $i_a \leq r_a$: closure under $\max$ with
another element $b = (r_b, i_b, 0, 0)$ gives
$\max(a, b) = (\max(r_a, r_b), \max(i_a, i_b), 0, 0)$, which can
have $i > r$ (take $r_a < r_b$ and $i_a > i_b$ with $i_a > r_b$).
So $D$ cannot be closed under both $\max$ and $\cdot_{\Hq}$ on $V$.
\end{proof}

\begin{remark}[The core tension]
The lattice and the monoid pull in opposite directions:
\begin{itemize}[nosep]
\item The \emph{lattice} wants a positively-ordered bounded domain where
  $\min$ and $\max$ are closed.
\item The \emph{Hamilton monoid} wants a norm-bounded domain where the
  signed cross terms can roam freely.
\end{itemize}
The component-wise order bounds each component independently; the
Hamilton product couples them through signed cross terms. The
incompatibility is structural: it arises from the \emph{minus signs}
in the quaternion multiplication table ($i^2 = j^2 = k^2 = -1$),
which are the same signs that produce non-commutativity. Removing them
would make the product commutative (and component-wise), recovering the
existing G-lattice.
\end{remark}

We now strengthen these negative results by ruling out four natural
candidate orders on~$B^4$.

\begin{theorem}[Cone impossibility]
\label{thm:cone-impossibility}
There is no proper cone $C \subsetneq \Hq \setminus \{0\}$ that is
invariant under both left and right multiplication by unit quaternions.
Consequently, no bi-invariant partial order on $\Hq$ can be defined via
a positive cone.
\end{theorem}

\begin{proof}
The group $\SU$ acts on $\Hq$ by conjugation: $q \mapsto uq\bar{u}$
for $|u| = 1$. This action preserves $\Real(q)$ and acts as $\mathrm{SO}(3)$
on $\Imag(q)$. Since $\mathrm{SO}(3)$ acts transitively on $S^2$,
the conjugation orbit of any pure imaginary unit quaternion is all of
$S^2 \subset \Imag(\Hq)$.

Suppose $C$ is a cone invariant under conjugation (a necessary
condition for bi-invariance). If $C$ contains any $q$ with
$\Imag(q) \neq 0$, then $C$ contains the entire conjugation orbit of
$q$, which includes elements with imaginary part pointing in every
direction of $S^2$. In particular, $C$ contains both $q'$ and elements
with opposite imaginary part, so $C$ contains elements in every open
half-space of $\Hq$. Thus $C$ cannot be a proper cone (a proper cone
is contained in a closed half-space).

For left-invariance specifically: if $q \in C$ and $|u| = 1$, then
$uq \in C$. But $\{uq : u \in S^3\}$ sweeps out the sphere of radius
$|q|$ in $\Hq$, so $C$ contains the full sphere $S^3_{|q|}$, which is
not contained in any proper cone.
\end{proof}

\begin{theorem}[$Q_8$ order impossibility]
\label{thm:q8-impossibility}
No non-trivial bi-invariant partial order exists on the quaternion
group $Q_8 = \{\pm 1, \pm i, \pm j, \pm k\}$.
\end{theorem}

\begin{proof}
Suppose $\leq$ is a partial order on $Q_8$ such that $a \leq b$
implies $ca \leq cb$ and $ac \leq bc$ for all $c \in Q_8$. Take any
$a \neq b$ with $a < b$. Left-multiplying by $c$ gives $ca < cb$;
right-multiplying gives $ac < bc$. We generate the \emph{order closure}:
the smallest partial order containing $a < b$ and closed under left
and right multiplication.

We claim this closure always produces a cycle.

\emph{Explicit example}: seed $1 < i$. Left-multiply by $i$:
$i < i^2 = -1$. Left-multiply $i < -1$ by $i$: $-1 < i \cdot (-1) = -i$.
Left-multiply $-1 < -i$ by $i$: $-i < i \cdot (-i) = 1$. The closure
now contains $1 < i < -1 < -i < 1$, a cycle contradicting
antisymmetry.

\emph{General argument}: for any $a \neq b$ with $a < b$,
left-multiply by $c = ba^{-1}$ (a unit): $ca < cb$, i.e.,
$b < cb$. The map $x \mapsto cx$ is a permutation of $Q_8$ of order
dividing~4 (every non-identity element of $Q_8$ has order~4). So
iterating gives $b < cb < c^2 b < c^3 b < c^4 b = b$, a
contradictory cycle. Exhaustive computation confirms that all $56$
ordered pairs $(a, b)$ with $a \neq b$ in $Q_8$ produce cycles within
4~steps.
\end{proof}

\begin{proposition}[Real-dominant subcone non-closure]
\label{prop:subcone}
For $\tau \in (0, 1)$, define the real-dominant subcone
$V_\tau = \{q \in B^4 : \Real(q) \geq \tau \cdot |q|\}$.
Then $V_\tau$ is not closed under the Hamilton product.
\end{proposition}

\begin{proof}
Take $q_1, q_2 \in S^3 \cap V_\tau$ with $\Real(q_i) = \tau$ and
imaginary parts aligned: $\Imag(q_1) = \Imag(q_2) = v$ with
$|v| = \sqrt{1 - \tau^2}$. Then:
\[
\Real(q_1 q_2) = \tau^2 - |\Imag(q_1)||\Imag(q_2)|\cos 0
  = \tau^2 - (1 - \tau^2) = 2\tau^2 - 1
\]
For $q_1 q_2 \in V_\tau$ we need $\Real(q_1 q_2) \geq \tau |q_1 q_2|
= \tau$, so we need $2\tau^2 - 1 \geq \tau$, equivalently
$2\tau^2 - \tau - 1 \geq 0$, which factors as $(2\tau + 1)(\tau - 1)
\geq 0$. Since $\tau > 0$, this holds only for $\tau \geq 1$. Thus
$V_\tau$ is not closed for any $\tau < 1$, and $V_1$ is the real line.
\end{proof}

\begin{proposition}[Norm-order collapse]
\label{prop:norm-collapse}
Define the \emph{norm order} on $B^4$ by $a \preceq b$ iff
$|a| \leq |b|$. Then $(B^4, \preceq, \cdot)$ satisfies:
\begin{enumerate}[nosep]
\item Monotonicity: $a \preceq b$ implies $c \cdot a \preceq c \cdot b$
  and $a \cdot c \preceq b \cdot c$.
\item The quotient $B^4 / {\sim}$ (where $a \sim b$ iff $|a| = |b|$)
  is isomorphic to $([0,1], \leq, \cdot_{\mathrm{prod}})$, the
  product BL-algebra on $[0,1]$.
\end{enumerate}
Consequently, the norm order erases all quaternionic structure: the
non-commutativity, axis coupling, and four-dimensionality are lost in
the quotient.
\end{proposition}

\begin{proof}
(1) follows from $|c \cdot a| = |c||a| \leq |c||b| = |c \cdot b|$.

(2) The map $\pi \colon B^4 \to [0,1]$, $q \mapsto |q|$ is a surjective
monoid homomorphism: $\pi(q_1 q_2) = |q_1 q_2| = |q_1||q_2|
= \pi(q_1) \cdot \pi(q_2)$, and $\pi(\mathbf{1}) = 1$. It is
order-preserving by definition of $\preceq$. The fibers $\pi^{-1}(r)$
are the 3-spheres $r \cdot S^3$ (for $r > 0$), on which $\preceq$
cannot distinguish elements. Thus $B^4/{\sim} \cong [0,1]$ as an
ordered monoid, and the product $\cdot_{\mathrm{prod}}$ on $[0,1]$
has the standard residuum $a \to b = \min(1, b/a)$.
\end{proof}

\subsection{Sphere exclusion}

The integrality condition (monoid identity = lattice top) imposes a
further constraint.

\begin{proposition}[Sphere exclusion]
\label{prop:sphere-exclusion}
Let $(D, \leq, \cdot)$ be an integral residuated lattice with
$D \subseteq B^4$, monoid identity $\mathbf{1} = (1,0,0,0)$, and
$\mathbf{1}$ as lattice top. If $a \in D \cap S^3$ and
$a^{-1} \in D$, then $a = \mathbf{1}$.
\end{proposition}

\begin{proof}
Since $\mathbf{1}$ is the lattice top, $a \leq \mathbf{1}$ for all
$a \in D$. Left-multiply by $a^{-1}$:
$a^{-1} \cdot a \leq a^{-1} \cdot \mathbf{1}$, i.e.,
$\mathbf{1} \leq a^{-1}$. But $a^{-1} \leq \mathbf{1}$ (top).
By antisymmetry, $a^{-1} = \mathbf{1}$, hence $a = \mathbf{1}$.
\end{proof}

\begin{corollary}
\label{cor:sphere-exclusion}
In any integral residuated lattice on $D \subseteq B^4$ with Hamilton
product, if $D$ is closed under inversion on $S^3$, then
$D \cap S^3 = \{\mathbf{1}\}$. The domain can contain at most the
identity on the unit sphere; all other elements must have
$|q| < 1$.
\end{corollary}

\subsection{The quaternionic trilemma}

The six results above---cone impossibility
(Theorem~\ref{thm:cone-impossibility}), $Q_8$ impossibility
(Theorem~\ref{thm:q8-impossibility}), subcone non-closure
(Proposition~\ref{prop:subcone}), norm-order collapse
(Proposition~\ref{prop:norm-collapse}), component-wise incompatibility
(Theorem~\ref{thm:order-incompatibility}), and sphere exclusion
(Proposition~\ref{prop:sphere-exclusion})---eliminate every natural
family of lattice orders on $B^4$.

\begin{theorem}[Quaternionic trilemma]
\label{thm:trilemma}
Let $T = \cdot_{\Hq}$ be the Hamilton product. Consider the three
properties: \textup{(a)}~associativity, \textup{(b)}~non-commutativity
arising from quaternionic cross-term coupling,
\textup{(c)}~compatible lattice order.
\begin{enumerate}[label=(\Roman*),nosep]
\item \textup{($V$-domains.)} For any $D \subseteq V$ containing
  $\mathbf{0}$ and $\mathbf{1}$ and closed under~$T$, at most two
  of \textup{(a)--(c)} hold simultaneously
  \textup{(}Theorem~\ref{thm:trilemma-V}\textup{)}.
\item \textup{($B^4$, equal norms.)} For orthogonal unit purely
  imaginary $\hat{u}, \hat{v}$ and $r \in (0,1)$, the sub-monoid
  $\langle r\hat{u}, r\hat{v}\rangle \cup \{\mathbf{0}, \mathbf{1}\}$
  satisfies \textup{(a)} and \textup{(b)} but admits no compatible
  lattice order \textup{(}Theorem~\ref{thm:trilemma-B4}\textup{)}.
\end{enumerate}
The equal-norm condition in \textup{(II)} is sharp: when $|a| \neq |b|$,
compatible monotone lattice orders exist
\textup{(}Remark~\ref{rem:sharpness}\textup{)}.
\end{theorem}

\begin{proof}
(I) is Theorem~\ref{thm:trilemma-V}; (II) is
Theorem~\ref{thm:trilemma-B4}. Sharpness is established by
computational verification in Remark~\ref{rem:sharpness}.
\end{proof}

The six elimination results complement the direct proofs:

\begin{theorem}[Elimination of natural orders]
\label{thm:elimination}
For the Hamilton monoid $(B^4, \cdot, \mathbf{1})$, properties
(a) and (b) hold, and the following families of orders fail to
provide~(c):
\begin{enumerate}[nosep]
\item Positive-cone orders: no proper bi-invariant cone exists
  (Theorem~\ref{thm:cone-impossibility}).
\item Finite subgroup orders: $Q_8$ admits no bi-invariant partial
  order (Theorem~\ref{thm:q8-impossibility}).
\item Subcone orders: no real-dominant subcone $V_\tau$ with
  $\tau < 1$ is closed (Proposition~\ref{prop:subcone}).
\item Norm-based orders: the norm order collapses to $[0,1]$, erasing
  non-commutativity (Proposition~\ref{prop:norm-collapse}).
\item Component-wise order: incompatible with Hamilton closure
  (Theorem~\ref{thm:order-incompatibility}).
\item Inversion-closed orders: the domain can meet $S^3$ only at
  $\{\mathbf{1}\}$ (Proposition~\ref{prop:sphere-exclusion}).
\end{enumerate}
\end{theorem}

\begin{proof}
Each item is proved in the cited result. The theorem collects them
to show that six independent approaches to constructing~(c) all fail
for the Hamilton monoid.
\end{proof}

\begin{remark}[Achievable pairs]
\label{rem:pairs}
Properties (a)$+$(b) are achieved by the Hamilton product on $B^4$.
Properties (a)$+$(c) are achieved by the G-lattice of~\citet{p11}
(component-wise min with the standard order on $V$). Properties
(b)$+$(c) are achieved by the Hamilton product on sub-monoids
$\langle r_1 \hat{u}, r_2 \hat{v}\rangle \cup \{\mathbf{0},
\mathbf{1}\}$ with $r_1 \neq r_2$: the multiplication is
non-commutative and a compatible monotone lattice order exists
(Remark~\ref{rem:sharpness}). Thus every pair of properties is
individually achievable; only the triple is impossible.
\end{remark}

\begin{remark}[Sharpness of the equal-norm condition]
\label{rem:sharpness}
The equal-norm hypothesis $|a| = |b|$ in
Theorem~\ref{thm:trilemma-B4} is essential, not merely a simplifying
assumption. When $|a| = |b| = r$, the identity $a^2 = b^2 = -r^2$
collapses four candidate level-2 products to three ($a^2, ab, ba$),
and the substitution $a \cdot b^2 = a \cdot a^2 = a^3$ creates a
cyclic swap among level-3 elements that forces a contradiction.

When $|a| \neq |b|$, $a^2 \neq b^2$ and level~2 has four distinct
elements. The chain $a^3 \leq a^2 b \leq a b^2$ no longer closes to
a cycle, because $a b^2 \neq a^3$. Computational verification
confirms that the sub-monoid
$\langle r_1 \hat{u}, r_2 \hat{v}\rangle \cup \{\mathbf{0},
\mathbf{1}\}$ with $r_1 \neq r_2$ admits a compatible monotone
lattice order at every tested depth (up to depth~5, with 18
elements). The solver finds satisfying assignments in very few
backtracking nodes, indicating that the lattice order is not
marginally feasible but genuinely unconstrained.
\end{remark}

%======================================================================
\section{The Alpha Bridge}
\label{sec:alpha}
%======================================================================

The tension between the commutative lattice (closed, monotone, no
coupling) and the Hamilton product (open, non-monotone, full coupling)
suggests a parameterized interpolation.

\begin{definition}[Alpha-parameterized product]
\label{def:alpha}
For $\alpha \in [0,1]$ and $q_1, q_2 \in V$, define:
\begin{equation}
\label{eq:alpha}
(q_1 \cdot_\alpha q_2) = \begin{pmatrix}
  r_1 r_2 - \alpha(i_1 i_2 + j_1 j_2 + k_1 k_2) \\
  r_1 i_2 + i_1 r_2 + \alpha(j_1 k_2 - k_1 j_2) \\
  r_1 j_2 + j_1 r_2 + \alpha(k_1 i_2 - i_1 k_2) \\
  r_1 k_2 + k_1 r_2 + \alpha(i_1 j_2 - j_1 i_2)
\end{pmatrix}
\end{equation}
At $\alpha = 0$, the cross terms vanish and the product becomes:
$(q_1 \cdot_0 q_2)_r = r_1 r_2$ and
$(q_1 \cdot_0 q_2)_m = r_1 q_{2,m} + q_{1,m} r_2$ for
$m \in \{i, j, k\}$. This is the \emph{diagonal part} of the Hamilton
product---the real component multiplies while the imaginary components
combine linearly, weighted by the real parts. It is \emph{not} the
component-wise (Hadamard) product $q_{1,m} q_{2,m}$, nor the
component-wise min used in the G-lattice of~\citet{p11}.
At $\alpha = 1$, the product is the full Hamilton
product~\eqref{eq:hamilton}.
\end{definition}

\begin{proposition}[Properties of the $\alpha$-family]
\label{prop:alpha}
\leavevmode
\begin{enumerate}[nosep]
\item \textbf{Identity}: $q \cdot_\alpha \mathbf{1} = \mathbf{1} \cdot_\alpha q = q$
  for all $\alpha \in [0,1]$.
\item \textbf{Annihilator}: $q \cdot_\alpha \mathbf{0} = \mathbf{0} \cdot_\alpha q
  = \mathbf{0}$ for all $\alpha$.
\item \textbf{Associativity}: $\cdot_\alpha$ is associative if and only
  if $\alpha \in \{0, 1\}$.
\item \textbf{Commutativity}: $\cdot_\alpha$ is commutative if and only
  if $\alpha = 0$.
\item \textbf{Specialization}: Under any restriction that sets two or
  more imaginary axes to zero, $\cdot_\alpha$ reduces to the same
  operation for all $\alpha$ (the cross terms vanish).
\end{enumerate}
\end{proposition}

\begin{proof}
(1)--(2) follow from direct substitution: the cross terms vanish when
one factor is $\mathbf{1}$ or $\mathbf{0}$.

(3)~At $\alpha = 0$: direct computation gives
$((q_1 \cdot_0 q_2) \cdot_0 q_3)_r = (r_1 r_2) r_3 = r_1 (r_2 r_3)
= (q_1 \cdot_0 (q_2 \cdot_0 q_3))_r$ and, for each imaginary axis~$m$,
$((q_1 \cdot_0 q_2) \cdot_0 q_3)_m
= (r_1 r_2) q_{3,m} + (r_1 q_{2,m} + q_{1,m} r_2) r_3
= r_1 r_2 q_{3,m} + r_1 q_{2,m} r_3 + q_{1,m} r_2 r_3
= r_1 (r_2 q_{3,m} + q_{2,m} r_3) + q_{1,m} (r_2 r_3)
= (q_1 \cdot_0 (q_2 \cdot_0 q_3))_m$. At $\alpha = 1$: associativity
is inherited from~$\Hq$.

For $0 < \alpha < 1$: the $\alpha$-product is a linear interpolation
between two different associative operations. Linear interpolation of
associative operations is generically non-associative. To verify: take
$a = (0, 0.5, 0, 0)$, $b = (0, 0, 0.5, 0)$, $c = (0.5, 0, 0, 0.5)$
and compute $(a \cdot_\alpha b) \cdot_\alpha c$ vs.\
$a \cdot_\alpha (b \cdot_\alpha c)$ for $\alpha = 0.5$. The
intermediate products generate terms in $\alpha^2$ that do not cancel.

(4)~The cross terms $j_1 k_2 - k_1 j_2$, etc., are antisymmetric in
$(q_1, q_2)$. For $\alpha > 0$, they contribute non-symmetrically; for
$\alpha = 0$, they vanish.

(5)~With two or more imaginary axes zeroed out, every cross term in
\eqref{eq:alpha} contains a factor from a zeroed axis and vanishes.
\end{proof}

\begin{remark}[Interpretation]
The parameter $\alpha$ controls the degree of \textbf{epistemic
coupling} between the four axes:
\begin{itemize}[nosep]
\item $\alpha = 0$: knowing ($r$), imagining ($i$), verifying ($j$),
  and selecting ($k$) are independent processes. The conjunction is
  component-wise.
\item $\alpha = 1$: the four processes interact through the full
  quaternionic structure. Imagining and verifying generate selection
  ($ij = k$). The order of operations matters.
\end{itemize}
The bridge is continuous but not algebraically smooth: associativity
holds only at the two endpoints. The loss of associativity for
intermediate $\alpha$ means that a gradual transition from commutative
to non-commutative logic is not algebraically coherent---the jump must
be discrete.
\end{remark}

%======================================================================
\section{Discussion}
\label{sec:discussion}
%======================================================================

\subsection{Summary of the landscape}

Table~\ref{tab:landscape} summarizes which properties are satisfied by
each construction.

\begin{table}[H]
\centering
\caption{Properties of candidate constructions. The trilemma manifests
as the impossibility of all seven checkmarks in any single column.}
\label{tab:landscape}
\small
\begin{tabularx}{\textwidth}{@{}l *{6}{>{\centering\arraybackslash}X} @{}}
\toprule
\textbf{Property}
  & \textbf{Min/$V$}
  & \textbf{Ham./$V$}
  & \textbf{Ham./$B^4$}
  & \textbf{Norm-ord.}
  & \textbf{$\alpha{=}0$}
  & \textbf{$\alpha{\in}(0,1)$} \\
\midrule
Identity      & $\checkmark$ & $\checkmark$ & $\checkmark$ & $\checkmark$ & $\checkmark$ & $\checkmark$ \\
Annihilator   & $\checkmark$ & $\checkmark$ & $\checkmark$ & $\checkmark$ & $\checkmark$ & $\checkmark$ \\
Associative   & $\checkmark$ & $\checkmark$ & $\checkmark$ & $\checkmark$ & $\checkmark$ & $\times$     \\
Closed        & $\checkmark$ & $\times$     & $\checkmark$ & $\checkmark$ & $\checkmark$ & partial      \\
Monotone      & $\checkmark$ & $\times$     & $\times$     & $\checkmark$ & $\checkmark$ & $\times$     \\
Commutative   & $\checkmark$ & $\times$     & $\times$     & eff.\ $\checkmark$ & $\checkmark$ & $\times$ \\
Axis coupling & $\times$     & $\checkmark$ & $\checkmark$ & $\times$     & $\times$     & partial      \\
\bottomrule
\end{tabularx}
\end{table}

\noindent The column ``Norm-ord.'' refers to the norm-order RL of
Proposition~\ref{prop:norm-collapse}: the Hamilton product is formally
non-commutative, but the order quotients it to the commutative
$[0,1]$ product algebra (``effectively commutative'').

The table confirms the trilemma: every column with axis coupling
($\checkmark$ in the last row) lacks monotonicity, and every column
with both associativity and monotonicity lacks axis coupling.

\subsection{Comparison with known structures}

\begin{table}[H]
\centering
\caption{The Hamilton product in the taxonomy of monoidal operations.}
\label{tab:comparison}
\begin{tabularx}{\textwidth}{@{}lYY@{}}
\toprule
\textbf{Structure} & \textbf{Properties} & \textbf{Relation to this work} \\
\midrule
t-norm \citep{hajek1998} & Commutative, associative, monotone on
  $[0,1]$ & Min t-norm is the current G-lattice conjunction \\
psBL-algebra \citep{hajek2003} & Non-commutative, two residuals on
  $[0,1]$ & Target structure; our monoid satisfies the monoid axioms
  but not the lattice axioms \\
Quantale \citep{mulvey1986} & Complete lattice + associative monoid &
  The unit ball monoid is a candidate if a compatible lattice order is
  found \\
Lambek calculus \citep{lambek1958} & Non-commutative conjunction
  in linguistics & Shares the two-residual structure; our cross terms
  give it geometric content \\
Quaternionic fuzzy sets \citep{dai2023} & Min/max on
  quaternion-valued memberships & No t-norm, no residuum;
  compatible with but independent of this work \\
Lattice-ordered groups \citep{kopytov1994} & Lattice + group;
  non-abelian examples exist on free groups & $\Hq^\times$ cannot be
  lattice-ordered: Theorem~\ref{thm:cone-impossibility} blocks the
  positive cone. Our monoid is not a group ($B^4$ has no inverses for
  $|q| < 1$), so the l-group theory does not directly apply \\
\bottomrule
\end{tabularx}
\end{table}

\subsection{Scope of the obstruction}
\label{sec:scope}

Theorem~\ref{thm:elimination} eliminates six families of orders.
The relationship between these eliminations and the sharp boundary
of Theorem~\ref{thm:trilemma} is characterised in
Remark~\ref{rem:sharpness}: the trilemma holds precisely under the
equal-norm condition $|a| = |b|$. Three further observations
sharpen the scope.

\begin{proposition}[No finite closed NC domain]
\label{prop:finite-domain}
Let $D \subset \mathrm{int}(B^4) \cup \{\mathbf{0}, \mathbf{1}\}$
be finite and closed under Hamilton product. Then
$D \subseteq \{\mathbf{0}, \mathbf{1}\}$.
\end{proposition}

\begin{proof}
Suppose $d \in D \setminus \{\mathbf{0}, \mathbf{1}\}$.
Then $d \in \mathrm{int}(B^4)$, so $0 < |d| < 1$.  The sequence
$d, d^2, d^3, \ldots$ lies in $D$ by closure, and
$|d^n| = |d|^n$ (norm multiplicativity of~$\Hq$). Since
$0 < |d| < 1$, the norms $|d|^n$ are pairwise distinct, so the
elements $d^n$ are pairwise distinct---an infinite subset of the
finite set~$D$, a contradiction.
\end{proof}

Consequently, any finite domain containing non-commuting elements
$a, b$ with $|a|, |b| < 1$ is not closed under the Hamilton product,
so the operation $T \colon D \times D \to D$ is not well-defined.
Combined with the $Q_8$ impossibility
(Theorem~\ref{thm:q8-impossibility}), which handles the boundary
case $|q| = 1$: no finite domain supports the trilemma triple.

A much stronger result holds for the natural domain
$V = [0,1] \times [-1,1]^3$ of the quaternionic logic.

\begin{theorem}[Angular accumulation]
\label{thm:angular}
Let $q \in V$ with $\Imag(q) \neq 0$. Then there exists a finite
$N$ such that $q^N \notin V$.
\end{theorem}

\begin{proof}
Write $q = |q|(\cos\theta + \sin\theta\,\hat{u})$ where
$\hat{u} = \Imag(q)/|\Imag(q)|$ is a unit pure imaginary quaternion
and $\theta = \arccos(\Real(q)/|q|) \in (0, \pi/2]$ (the bounds
follow from $\Real(q) \geq 0$ and $|\Imag(q)| > 0$). By
De Moivre's theorem for quaternions ($\hat{u}^2 = -1$ gives the
same exponential structure as complex numbers):
\[
  q^n = |q|^n\bigl(\cos(n\theta) + \sin(n\theta)\,\hat{u}\bigr).
\]
In particular, $\Real(q^n) = |q|^n \cos(n\theta)$.

Since $\theta > 0$, there exists a smallest $N$ with
$N\theta > \pi/2$; explicitly, $N \leq \lceil\pi/(2\theta)\rceil + 1$.
At $n = N$: if $\cos(N\theta) < 0$, then $\Real(q^N) < 0$, so
$q^N \notin V$. If $\cos(N\theta) = 0$ (i.e.\ $N\theta = \pi/2$
exactly), then $\cos((N{+}1)\theta) = -\sin\theta < 0$, so
$q^{N+1} \notin V$.
\end{proof}

\begin{corollary}[Only commutative sub-monoids of $V$]
\label{cor:commutative-submonoid}
Every sub-monoid of $(V, \cdot_{\Hq})$ contained in $V$ is
commutative. That is: if $D \subseteq V$ is closed under the
Hamilton product and contains $\mathbf{1}$, then $ab = ba$ for all
$a, b \in D$.
\end{corollary}

\begin{proof}
Suppose $d \in D$ has $\Imag(d) \neq 0$. By closure, $d^n \in D
\subseteq V$ for all $n \geq 1$. But Theorem~\ref{thm:angular}
gives $d^N \notin V$ for some finite $N$, a contradiction. So every
$d \in D$ is real: $D \subseteq \{(r,0,0,0) : r \in [0,1]\}$.
Real quaternions commute.
\end{proof}

\begin{theorem}[Quaternionic trilemma for $V$-domains]
\label{thm:trilemma-V}
Let $D \subseteq V$ contain $\mathbf{0}$ and $\mathbf{1}$, and let
$T \colon D \times D \to D$ be the Hamilton product restricted
to~$D$. Then at most two of \{(a)~associativity,
(b)~non-commutativity, (c)~compatible lattice order\} hold.
\end{theorem}

\begin{proof}
Condition (c) requires $T \colon D \times D \to D$ to be
well-defined, hence $D$ must be closed under $T$.
Corollary~\ref{cor:commutative-submonoid} then forces $D$ to be
commutative, contradicting~(b).
\end{proof}

The trilemma is therefore a \emph{theorem} on the natural domain
$V$ of the quaternionic logic. The mechanism is angular entropy
production: each non-real element carries an angle
$\theta = \arccos(\Real(q)/|q|) > 0$, and iterated multiplication
accumulates this angle linearly ($n\theta$). The boundary $\Real
\geq 0$ of $V$ can absorb at most $\pi/2$ radians of angular
entropy; beyond this threshold the system exits $V$. The
non-commutativity and the boundary violation share the same algebraic
root: the relation $\hat{u}^2 = -1$ for pure imaginary unit
quaternions.

A stronger result holds for the full ball $B^4$, where powers remain
in $B^4$ and the angular accumulation argument does not apply.
The obstruction shifts from closure failure to a lattice-theoretic
incompatibility within the non-commutative sub-monoid itself.

\begin{proposition}[Level-2 incomparability in $B^4$]
\label{prop:level2-incomparable}
Let $r \in (0,1)$, $a = ri$, $b = rj$, and let $\leq$ be any
partial order on $M = \langle a, b \rangle \cup \{\mathbf{0},
\mathbf{1}\}$ with $\mathbf{0}$ as bottom, $\mathbf{1}$ as top,
and the Hamilton product monotone on both sides. Then:
\begin{enumerate}[label=\textup{(\alph*)},nosep]
\item $a \| b$ \textup{(}the generators are incomparable\textup{)};
\item $a^2$, $ab$, and $ba$ are pairwise incomparable.
\end{enumerate}
\end{proposition}

\begin{proof}
\textbf{(a)} Suppose $a \leq b$. Left-multiply by $a$:
$a^2 \leq ab$, i.e.\ $-r^2 \leq r^2k$. Right-multiply by $b$:
$ab \leq b^2$, i.e.\ $r^2k \leq -r^2$. Together:
$r^2k = -r^2$, which is false. The case $b \leq a$ is symmetric
(yielding $-r^2k = -r^2$).

\textbf{(b)} We show the case $a^2 \leq ab$ in detail; the
remaining five are analogous.

Suppose $a^2 \leq ab$. Left-multiply by $a$: $a^3 \leq a^2 b$,
i.e.\ $-r^3 i \leq -r^3 j$. Right-multiply by $b$:
$a^2 b \leq ab^2 = a^3$, i.e.\ $-r^3 j \leq -r^3 i$.
By antisymmetry, $-r^3 i = -r^3 j$, contradicting $i \neq j$.
In each of the six cases, two applications of monotonicity
(one left, one right multiplication by generators) produce a cycle
$p \leq q \leq p$ between distinct level-3 quaternions. The full
case analysis is verified computationally in the supplementary
script \texttt{b4\_proof\_verify.py}.
\end{proof}

\begin{theorem}[Lattice obstruction for $B^4$]
\label{thm:trilemma-B4}
Let $\hat{u}, \hat{v}$ be orthogonal unit purely imaginary
quaternions and $r \in (0,1)$. The sub-monoid
$M = \langle r\hat{u},\, r\hat{v} \rangle
\cup \{\mathbf{0},\,\mathbf{1}\}$
of $(B^4, \cdot_{\Hq})$ admits no partial order $\leq$ satisfying
simultaneously:
\begin{enumerate}[label=\textup{(\roman*)},nosep]
\item $\mathbf{0}$ is bottom and $\mathbf{1}$ is top;
\item $x \leq y$ implies $zx \leq zy$ and $xz \leq yz$ for all
  $z \in M$;
\item $(M, \leq)$ is a lattice.
\end{enumerate}
\end{theorem}

\begin{proof}
By $\SU$ conjugation (which preserves norms, products, and the
order structure), we may assume $\hat{u} = i$, $\hat{v} = j$,
so $a = ri$ and $b = rj$. The products $a^2 = b^2 = -r^2$,
$ab = r^2 k$, $ba = -r^2 k$ are pairwise distinct.

\emph{Forced bounds.} From $0 \leq a, b \leq 1$ and (ii):
for any word $w = w_1 w_2 \cdots w_n$ with $w_i \in \{a, b\}$ and
$n \geq 2$, induction on $n$ using $w_n \leq 1$ gives
$w \leq w_1 w_2$. In particular, every element of $M$ at level
$\geq 3$ is below its two-generator prefix in $\{a^2, ab, ba\}$.

\emph{Upper bounds of $\{a^2, ba\}$.} We claim that the only
elements of $M$ above both $a^2$ and $ba$ are $\{a, b, 1\}$.

Level~1: $a \geq a^2$ (from $a \leq 1$, multiply by $a$) and
$a \geq ba$ (from $b \leq 1$, right-multiply by $a$). Similarly
$b \geq a^2$ (since $a^2 = b^2 \leq b$) and $b \geq ba$
(from $a \leq 1$, left-multiply by $b$).

Level~2: any of $ab \geq a^2$, $ba \geq a^2$, or $a^2 \geq ba$
contradicts Proposition~\ref{prop:level2-incomparable}(b).

Level~$\geq 3$: if $w \geq a^2$ for a word of length $n \geq 3$,
then $a^2 \leq w \leq w_1 w_2$, forcing $a^2 \leq w_1 w_2$.
Proposition~\ref{prop:level2-incomparable}(b) excludes
$w_1 w_2 \in \{ab, ba\}$, so $w_1 w_2 = a^2$, giving
$w = a^2$ by antisymmetry. But $|w| = r^n < r^2 = |a^2|$ for
$n \geq 3$, so $w \neq a^2$---a contradiction.

$\mathbf{0}$: if $0 \geq a^2$ then $a^2 = 0$, contradicting
$a^2 = -r^2 \neq 0$.

\emph{No join.} The upper bounds of $\{a^2, ba\}$ in $M$ are
exactly $\{a, b, 1\}$. Both $a$ and $b$ are minimal among these
($a < 1$, $b < 1$, but $a \| b$ by
Proposition~\ref{prop:level2-incomparable}(a)). Two incomparable
minimal upper bounds admit no least upper bound. The join
$a^2 \vee ba$ does not exist, so $(M, \leq)$ is not a lattice.
\end{proof}

The obstruction mechanism differs fundamentally from the $V$-domain
case. In $V$, non-commutativity is eliminated entirely (every
sub-monoid is commutative). In $B^4$, non-commutative sub-monoids
exist---the Hamilton product is closed on $B^4$---but they cannot
carry a compatible lattice order. The root cause is the same in
both cases: the relation $\hat{u}^2 = -1$ for imaginary unit
quaternions.

The equal-norm hypothesis is sharp. When $|a| \neq |b|$, the
collapse $a^2 = b^2$ that drives the cyclic swap at level~3 no
longer holds, and compatible monotone lattice orders exist
(Remark~\ref{rem:sharpness}). This delineates the precise
boundary of the trilemma within $B^4$: non-commutativity and
compatible order coexist exactly when the generators have unequal
norms, restoring sufficient algebraic room for the lattice.

\begin{remark}[Quaternionic specificity]
\label{rem:specificity}
The obstruction is specific to the quaternionic setting. The only
normed division algebras over $\mathbb{R}$ are $\mathbb{R}$,
$\mathbb{C}$, $\mathbb{H}$, and $\mathbb{O}$ (Hurwitz's theorem).
Among these: $\mathbb{R}$ and $\mathbb{C}$ are commutative
(no property~(b)), $\mathbb{O}$ is non-associative (no
property~(a)), and $\mathbb{H}$ is the subject of this paper.
Clifford algebras $\mathrm{Cl}(p,q)$ with $p + q \geq 3$ are not
division algebras and thus admit zero divisors, precluding
well-defined residuals. The trilemma is therefore a statement about
the \emph{unique} algebraic setting---$\mathbb{H}$---where all
three properties can be meaningfully posed.
\end{remark}

\begin{remark}[Residuation under the norm order]
\label{rem:residuation-norm}
One might ask whether the division-based residuals
(Definition~\ref{def:implications}), truncated to $B^4$ via
$a \to_L b = \pi_{B^4}(a^{-1} b)$, preserve the adjunction
approximately. Proposition~\ref{prop:norm-collapse} settles this:
the only monotone order on $B^4$ compatible with the Hamilton
product is the norm order $|q_1| \leq |q_2|$, which collapses the
lattice to $[0,1]$ with the product t-norm. Under this order, the
residual is $|a| \to |b| = \min(1, |b|/|a|)$ and the adjunction
holds \emph{exactly}---but the quaternionic structure beyond the
norm is invisible. No intermediate approximation between full
quaternionic residuation and norm collapse is available, because the
monotonicity failures (Theorem~\ref{thm:monotonicity-failure}) are
structural, not perturbative.
\end{remark}

%======================================================================
\section{Conclusion}
\label{sec:conclusion}
%======================================================================

The Hamilton product of quaternions is a natural candidate for a
non-commutative conjunction in quaternionic logic: it couples the four
epistemic axes, satisfies bilateral identity and associativity, admits
left and right implications with modus ponens, and reduces to standard
commutative t-norms under every proper restriction. Its
non-commutativity is emergent---a genuinely four-dimensional phenomenon
absent from all lower-dimensional projections.

The obstructions to a full residuated lattice are equally deep. Six
independent impossibility results eliminate every natural family of
lattice orders on $B^4$: bi-invariant cones (SU(2) transitivity),
finite group orders ($Q_8$ exhaustive), real-dominant subcones
(analytical bound), norm-based orders (1D collapse), component-wise
orders (closure failure), and inversion-closed domains (sphere
exclusion). These impossibilities hold in the equal-norm regime.
For unequal-norm generators, the obstruction dissolves: every pair
of properties is individually achievable (Remark~\ref{rem:pairs}),
and the triple is impossible precisely when $|a| = |b|$.

This elimination is not a deficiency of any particular
construction---it traces to the defining relations
$i^2 = j^2 = k^2 = -1$. The minus signs that produce
non-commutativity are the same signs that break monotonicity---and,
via De Moivre's theorem, the same signs that force iterated products
out of the truth-value domain~$V$.

The quaternionic trilemma is a theorem in two settings. For
$V$-domains (Theorem~\ref{thm:trilemma-V}), the proof is that every
sub-monoid of $V$ under the Hamilton product is commutative
(Corollary~\ref{cor:commutative-submonoid}), because the angular
entropy $\theta = \arccos(\Real(q)/|q|)$ of any non-real element
accumulates linearly under iteration, exceeding $V$'s capacity
$\pi/2$ in finitely many steps (Theorem~\ref{thm:angular}). For
equal-norm sub-monoids of $B^4$
(Theorem~\ref{thm:trilemma-B4}), the proof is that the collapse
$a^2 = b^2$ forces pairwise incomparability of level-2 products
and prevents the existence of joins. The equal-norm condition is
sharp: unequal-norm sub-monoids admit compatible lattice orders
(Remark~\ref{rem:sharpness}), so the trilemma boundary is
precisely the locus $|a| = |b|$. Six family eliminations
(Theorem~\ref{thm:elimination}) and nine computationally surveyed
candidates (Table~\ref{tab:landscape}) independently corroborate
the impossibility in the equal-norm regime.

%======================================================================
% BIBLIOGRAPHY
%======================================================================

\bibliographystyle{plainnat}

\begin{thebibliography}{99}

\bibitem[Ciungu(2006)]{ciungu2006}
Ciungu, L.C. (2006).
Classes of residuated lattices.
\textit{Annals of the University of Craiova, Mathematics and Computer
Science Series}, 33, 56--70.

\bibitem[Dai(2023)]{dai2023}
Dai, S. (2023).
Quaternionic fuzzy sets.
\textit{Axioms}, 12(5), 490.

\bibitem[Esteva \& Godo(2001)]{esteva2001}
Esteva, F. \& Godo, L. (2001).
Monoidal t-norm based logic: towards a logic for left-continuous t-norms.
\textit{Fuzzy Sets and Systems}, 124(3), 271--288.

\bibitem[Galatos et~al.(2007)]{galatos2007}
Galatos, N., Jipsen, P., Kowalski, T. \& Ono, H. (2007).
\textit{Residuated Lattices: An Algebraic Glimpse at Substructural Logics}.
Elsevier.

\bibitem[H\'ajek(1998)]{hajek1998}
H\'ajek, P. (1998).
\textit{Metamathematics of Fuzzy Logic}.
Kluwer.

\bibitem[H\'ajek(2003)]{hajek2003}
H\'ajek, P. (2003).
Observations on non-commutative fuzzy logic.
\textit{Soft Computing}, 7(1), 38--43.

\bibitem[Kopytov \& Medvedev(1994)]{kopytov1994}
Kopytov, V.M. \& Medvedev, N.Ya. (1994).
\textit{The Theory of Lattice-Ordered Groups}.
Kluwer.

\bibitem[Lambek(1958)]{lambek1958}
Lambek, J. (1958).
The mathematics of sentence structure.
\textit{American Mathematical Monthly}, 65(3), 154--170.

\bibitem[Mulvey(1986)]{mulvey1986}
Mulvey, C.J. (1986).
\&.
\textit{Rendiconti del Circolo Matematico di Palermo, Serie II,
Supplemento}, 12, 99--104.

\bibitem[Ornelas Brand(2026a)]{p11}
Ornelas Brand, J.A. (2026).
Quaternionic Logic: A G-Lattice Unifying Boolean and Fuzzy Frameworks.
Manuscript in preparation.

\bibitem[Restall(2000)]{restall2000}
Restall, G. (2000).
\textit{An Introduction to Substructural Logics}.
Routledge.

\end{thebibliography}

\end{document}
